\documentclass[dvipdfmx]{beamer}
%テーマ
\usetheme[numbering=counter,shape=pentagon]{LightTheme}
%パッケージ
\usepackage{pxjahyper}
\usepackage{tikz}
\usepackage[absolute,overlay]{textpos}
\usepackage{xcolor}
%tikzライブラリ
\usetikzlibrary{arrows.meta}
\usetikzlibrary{positioning}
\usetikzlibrary{shapes.callouts}
%フォント
\renewcommand{\kanjifamilydefault}{\gtdefault}

%タイトル設定
\title{TechRingにかける想い}
\subtitle{学科を超えて}
\author[Jin]{石田　優仁}
\institute[TechRing]{奈良高専TechRing}
\date[2023/04/30]{2023/04/30}

\begin{document}

\begin{frame}[plain,noframenumbering]
  \begin{flushright}
      \includegraphics[width=2.3cm]{pic/Logo.png}
  \end{flushright}
  \titlepage
\end{frame}

\begin{frame}{高専どんなところ？}

  \begin{block}{国立高等専門学校の特徴}
    \begin{itemize}
      \setlength{\itemsep}{1mm}
      \item 5年制、一貫のカリキュラム
      \item 電気電子系や機械系に限らず、物質や建築、商船もある
      \item 高専祭という文化祭での文化発表が活発
      \item ロボコンやデザコン、プロコンなどのコンテストが充実
      \item 高専カンファレンスや様々な技術コミュニティでの交流
    \end{itemize}
  \end{block}

  \begin{alertblock}{奈良高専の特徴}
      \begin{itemize}
        \setlength{\itemsep}{1mm}
        \item 来年で創設60周年を迎える
        \item 機械、電気、電子制御、情報、物質化学の5学科
        \item 最寄り駅からバスで20分ほど
        \item 女子枠推薦がある
      \end{itemize}
  \end{alertblock}

\end{frame}

\begin{frame}{高専スタートアップ事業について}{時代は大高専時代に...}

  \begin{alertblock}{高専スタートアップ事業概要}
      \begin{itemize}
        \setlength{\itemsep}{1mm}
        \item アントレプレナーシップ教育に取り組む\\
        国公私立高専の戦略的な取組に必要な経費を補助
        \item 高専生の活動を後押しし\\
        学生が自らの技術を用いた社会課題解決に取り組み
        \item 予算は57高専で合計\alert{60億円}も
      \end{itemize}
  \end{alertblock}

  \vspace{-3mm}

  \begin{center}
    \begin{tikzpicture}
      \draw[-{Latex[length=5mm]},line width = 3pt,red] (0,.7) -- (0,0);
    \end{tikzpicture}
  \end{center}

  \vspace{-4.5mm}

  \begin{greyblock}{高専スタートアップ事業STEP}
    \begin{enumerate}
      \setlength{\itemsep}{1mm}
      \item アントレプレナーシップや企業について知る
      \item 高専生のチャレンジを応援するための環境整備\\
      $\Rightarrow$\alert{起業家工房}の整備
      \item 地域社会の課題解決に挑戦
    \end{enumerate}
  \end{greyblock}
\end{frame}

\begin{frame}{TechRingとは？}

  \begin{description}[一つの輪に]
    \item[学科の壁] 各学科には様々な長所や短所が混在するが...\par
    各学科が壁で隔てられ、\\長所を活かすことも短所を補うこともできない！\par
    \hspace{0.3\textwidth}
    \begin{tikzpicture}
      \draw[-{Latex[length=5mm]},line width = 3pt,red] (0,.7) -- (0,0);
    \end{tikzpicture}
    \vspace{-1.5mm}
    \item[一つの輪に] 学科を超えた相互補完型の交流で\\\structure{究極のアントレプレナーシップ}を目指す！
  \end{description}

  \begin{columns}
    \begin{column}{0.55\textwidth}
      \begin{center}
        \begin{block}{TechRing ３つの起業}
          \vspace{2mm}
          \begin{enumerate}
            \setlength{\itemsep}{3mm}
            \item 学科交流イベントの起業
            \item 技術交流プロジェクトの起業
            \item 共同開発プロジェクトの起業
          \end{enumerate}
        \end{block}
      \end{center}
    \end{column}

    \begin{column}{0.40\textwidth}
      \begin{center}
        \begin{tikzpicture}[scale=1.4]
          %座標設定
          \coordinate (A) at (0,-.1);
          \coordinate (B) at (0,1);
          \coordinate (C) at ({cos(9*pi/10 r)},{sin(9*pi/10 r)});
          \coordinate (D) at ({cos(13*pi/10 r)},{sin(13*pi/10 r)});
          \coordinate (E) at ({cos(17*pi/10 r)},{sin(17*pi/10 r)});
          \coordinate (F) at ({cos(21*pi/10 r)},{sin(21*pi/10 r)});
          %線引き
          \draw[line width = 2pt] (F)--(B)--(C)--(D)--(E)--cycle;
          %nodeのスタイル指定
          \tikzset{block/.style={circle,white,fill=cyan,minimum size=1cm,inner sep=.5pt}};
          %node（テキストボックス的な）を指定座標に配置
          \node[inner sep=0pt] at (A) {\includegraphics[width=1.5cm]{pic/Logo.png}};
          \node[block] at (B) {\footnotesize MeCafe};
          \node[block] at (C) {\footnotesize メカ研};
          \node[block] at (D) {\footnotesize 電研};
          \node[block] at (E) {\footnotesize シス研};
          \node[block] at (F) {\footnotesize 情研};
        \end{tikzpicture}
      \end{center}
    \end{column}
  \end{columns}

\end{frame}

\begin{frame}{TechRing ３つの目的 その１}
  \begin{greyblock}{学科を超えた技術交流とその起業のために}
    \vspace{1mm}
    \begin{footnotesize}
      \begin{itemize}
        \item \alert{学内技術イベントを計画・運営}
        \item \alert{学内ロボットコンテストの共同開催}
        \item \alert{高専祭学科展にて相互連携}
        \item TechRing内技術交流
        \item 他高専の有志・団体との交流
        \item 各種コンテスト、展示会への出場
        \item 共同開発プロジェクトの企画
        \item チャレンジプロジェクトに参加
      \end{itemize}
    \end{footnotesize}
  \end{greyblock}

  \begin{columns}
    \begin{column}{0.30\textwidth}
      \begin{alertblock}{学内イベント}
        \begin{footnotesize}
          TechRing ３つの起業
          \begin{enumerate}
            \item 学科交流\\Eventの起業
          \end{enumerate}
        \end{footnotesize}
      \end{alertblock}
    \end{column}

    \begin{column}{0.30\textwidth}
      \begin{block}{技術交流}
        \begin{footnotesize}
          TechRing ３つの起業
          \begin{enumerate}
            \setcounter{enumi}{1}
            \item 技術交流\\Projectの起業
          \end{enumerate}
        \end{footnotesize}
      \end{block}
    \end{column}

    \begin{column}{0.30\textwidth}
      \begin{block}{共同制作}
        \begin{footnotesize}
          TechRing ３つの起業 \par
          \begin{enumerate}
            \setcounter{enumi}{2}
            \item 共同開発\\Projectの起業
          \end{enumerate}
        \end{footnotesize}
      \end{block}
    \end{column}
  \end{columns}
\end{frame}

\begin{frame}{TechRing ３つの活動 その１}{イベント運営}
  \begin{columns}
    \begin{column}{0.55\textwidth}
      \begin{alertblock}{TechCafe}
        \begin{itemize}
          \item 年に一度の学内一大\\LTイベントとして開催
          \item 技術発表に限らない\\幅広いテーマ設定
          \item 高専祭のアフターイベント\\としてバックアップ
        \end{itemize}
      \end{alertblock}

      \begin{block}{他にも}
        \begin{footnotesize}
          \begin{itemize}
            \item 高専祭学科展にて連携
            \item 定期的な学科交流イベント
            \item 学内ロボコンできたらいいな
          \end{itemize}
        \end{footnotesize}
      \end{block}
    \end{column}

    \begin{column}{0.4\textwidth}
      \includegraphics[width=4.5cm]{pic/TechCafe.png}
    \end{column}
  \end{columns}
\end{frame}

\begin{frame}{TechRing ３つの目的 その２}
  \begin{greyblock}{学科を超えた技術交流とその起業のために}
    \vspace{1mm}
    \begin{footnotesize}
      \begin{itemize}
        \item 学内技術イベントを計画・運営
        \item 学内ロボットコンテストの共同開催
        \item 高専祭学科展にて相互連携
        \item \alert{TechRing内技術交流}
        \item \alert{他高専の有志・団体との交流}
        \item \alert{各種コンテスト、展示会への出場}
        \item 共同開発プロジェクトの企画
        \item チャレンジプロジェクトに参加
      \end{itemize}
    \end{footnotesize}
  \end{greyblock}

  \begin{columns}
    \begin{column}{0.30\textwidth}
      \begin{block}{学内イベント}
        \begin{footnotesize}
          TechRing ３つの起業
          \begin{enumerate}
            \item 学科交流\\Eventの起業
          \end{enumerate}
        \end{footnotesize}
      \end{block}
    \end{column}

    \begin{column}{0.30\textwidth}
      \begin{alertblock}{技術交流}
        \begin{footnotesize}
          TechRing ３つの起業
          \begin{enumerate}
            \setcounter{enumi}{1}
            \item 技術交流\\Projectの起業
          \end{enumerate}
        \end{footnotesize}
      \end{alertblock}
    \end{column}

    \begin{column}{0.30\textwidth}
      \begin{block}{共同制作}
        \begin{footnotesize}
          TechRing ３つの起業 \par
          \begin{enumerate}
            \setcounter{enumi}{2}
            \item 共同開発\\Projectの起業
          \end{enumerate}
        \end{footnotesize}
      \end{block}
    \end{column}
  \end{columns}
\end{frame}

\begin{frame}{TechRing ３つの活動 その２}{技術交流}
  \begin{alertblock}{TechRingミーティング}
    \begin{itemize}
      \item 定期的な活動報告会
      \item 他団体からノウハウを借りる講演や実習の開催
    \end{itemize}
  \end{alertblock}

  \begin{alertblock}{Maker Faire Kyoto 2023 出展}
    \begin{itemize}
      \item 各団体が製作した製作品のアウトプット
      \item 共同展示によって連携し、奈良高専全体の広報にも
    \end{itemize}
  \end{alertblock}

  \begin{block}{製作品外部公開}
    \begin{itemize}
      \item Twitterを使用した発信
      \item 外部公開によって奈良高専外部との交流も深める
    \end{itemize}
  \end{block}
\end{frame}

\begin{frame}{TechRing ３つの目的 その３}
  \begin{greyblock}{学科を超えた技術交流とその起業のために}
    \vspace{1mm}
    \begin{footnotesize}
      \begin{itemize}
        \item 学内技術イベントを計画・運営
        \item 学内ロボットコンテストの共同開催
        \item 高専祭学科展にて相互連携
        \item TechRing内技術交流
        \item 他高専の有志・団体との交流
        \item \alert{各種コンテスト、展示会への出場}
        \item \alert{共同開発プロジェクトの企画}
        \item \alert{チャレンジプロジェクトに参加}
      \end{itemize}
    \end{footnotesize}
  \end{greyblock}

  \begin{columns}
    \begin{column}{0.30\textwidth}
      \begin{block}{学内イベント}
        \begin{footnotesize}
          TechRing ３つの起業
          \begin{enumerate}
            \item 学科交流\\Eventの起業
          \end{enumerate}
        \end{footnotesize}
      \end{block}
    \end{column}

    \begin{column}{0.30\textwidth}
      \begin{block}{技術交流}
        \begin{footnotesize}
          TechRing ３つの起業
          \begin{enumerate}
            \setcounter{enumi}{1}
            \item 技術交流\\Projectの起業
          \end{enumerate}
        \end{footnotesize}
      \end{block}
    \end{column}

    \begin{column}{0.30\textwidth}
      \begin{alertblock}{共同制作}
        \begin{footnotesize}
          TechRing ３つの起業 \par
          \begin{enumerate}
            \setcounter{enumi}{2}
            \item 共同開発\\Projectの起業
          \end{enumerate}
        \end{footnotesize}
      \end{alertblock}
    \end{column}
  \end{columns}
\end{frame}

\begin{frame}{TechRing ３つの活動 その３}{共同開発}
  \begin{columns}
    \begin{column}{0.63\textwidth}
      \begin{block}{開発環境の整備}
        \begin{itemize}
          \item 起業家工房の環境を整備
          \item 起業家工房による\\有志個人のものづくり支援も
        \end{itemize}
      \end{block}
    \end{column}

    \begin{column}{0.28\textwidth}
      \begin{center}
        \includegraphics[width=2.5cm]{pic/Team.png}
      \end{center}
    \end{column}
  \end{columns}

  \begin{alertblock}{共同開発プロジェクト}
    \begin{itemize}
      \item 学科を跨いだ人員を確保し、コンテスト共同参加も
      \item 共同で製作した製作品の展示、アウトプット
      \item 広く外部公開し、学外有志の意見も取り入れフィードバック
      \item 製作したモノの担当したパートを各団体に持ち帰って共有
    \end{itemize}
  \end{alertblock}
\end{frame}

\begin{frame}{TechRingの組成}{３つの起業を５つの団体で}
  \begin{block}{何処から誰を招集するか}
    \vspace{2mm}
    \begin{description}
      \setlength{\itemsep}{2mm}
      \item[招集基準] 各学科を代表する団体を招集したい
      \item[判断素材] 高専祭学科展の展示に携わった経験のある団体
      \item[招集結果] ４学科から５団体の参加が確定！
    \end{description}
  \end{block}

  \begin{center}
    \begin{tikzpicture}
      \draw[-{Latex[length=5mm]},line width = 3pt,red] (0,.7) -- (0,0);
    \end{tikzpicture}
  \end{center}

  \begin{Large}
    各学科各団体の\alert{長所を活かし} \vfill \hspace{5em}\alert{短所を補う}プロジェクトとして発足！
  \end{Large}
\end{frame}

\begin{frame}{TechRingを創る仲間たち その１}{MeCafe}
  機械研究会にフォーカス\\
  活動の紹介と出展の紹介を
  \begin{textblock*}{0.4\linewidth}(280pt,180pt)
    \begin{tikzpicture}[scale=1]
      %座標設定
      \coordinate (A) at (0,-.1);
      \coordinate (B) at (0,1);
      \coordinate (C) at ({cos(9*pi/10 r)},{sin(9*pi/10 r)});
      \coordinate (D) at ({cos(13*pi/10 r)},{sin(13*pi/10 r)});
      \coordinate (E) at ({cos(17*pi/10 r)},{sin(17*pi/10 r)});
      \coordinate (F) at ({cos(21*pi/10 r)},{sin(21*pi/10 r)});
      %線引き
      \draw[line width = 2pt] (F)--(B)--(C)--(D)--(E)--cycle;
      %nodeのスタイル指定
      \tikzset{block/.style={circle,white,fill=cyan,minimum size=.85cm,inner sep=0pt}};
      %node（テキストボックス的な）を指定座標に配置
      \node[inner sep=0pt] at (A) {\includegraphics[width=1cm]{pic/Logo.png}};
      \node[block,fill=LightUIRed] at (B) {\scriptsize MeCafe};
      \node[block] at (C) {\scriptsize メカ研};
      \node[block] at (D) {\scriptsize 電研};
      \node[block] at (E) {\scriptsize シス研};
      \node[block] at (F) {\scriptsize 情研};
    \end{tikzpicture}
  \end{textblock*}
\end{frame}

\begin{frame}{TechRingを創る仲間たち その２}{機械研究会}
  MeCafeにフォーカス\\
  活動の紹介と出展の紹介を
  \begin{textblock*}{0.4\linewidth}(280pt,180pt)
    \begin{tikzpicture}[scale=1]
      %座標設定
      \coordinate (A) at (0,-.1);
      \coordinate (B) at (0,1);
      \coordinate (C) at ({cos(9*pi/10 r)},{sin(9*pi/10 r)});
      \coordinate (D) at ({cos(13*pi/10 r)},{sin(13*pi/10 r)});
      \coordinate (E) at ({cos(17*pi/10 r)},{sin(17*pi/10 r)});
      \coordinate (F) at ({cos(21*pi/10 r)},{sin(21*pi/10 r)});
      %線引き
      \draw[line width = 2pt] (F)--(B)--(C)--(D)--(E)--cycle;
      %nodeのスタイル指定
      \tikzset{block/.style={circle,white,fill=cyan,minimum size=.85cm,inner sep=0pt}};
      %node（テキストボックス的な）を指定座標に配置
      \node[inner sep=0pt] at (A) {\includegraphics[width=1cm]{pic/Logo.png}};
      \node[block] at (B) {\scriptsize MeCafe};
      \node[block,fill=LightUIRed] at (C) {\scriptsize メカ研};
      \node[block] at (D) {\scriptsize 電研};
      \node[block] at (E) {\scriptsize シス研};
      \node[block] at (F) {\scriptsize 情研};
    \end{tikzpicture}
  \end{textblock*}
\end{frame}

\begin{frame}{TechRingを創る仲間たち その３}{電気技術研究会}
  電気技術研究会にフォーカス\\
  活動の紹介と出展の紹介を
  \begin{textblock*}{0.4\linewidth}(280pt,180pt)
    \begin{tikzpicture}[scale=1]
      %座標設定
      \coordinate (A) at (0,-.1);
      \coordinate (B) at (0,1);
      \coordinate (C) at ({cos(9*pi/10 r)},{sin(9*pi/10 r)});
      \coordinate (D) at ({cos(13*pi/10 r)},{sin(13*pi/10 r)});
      \coordinate (E) at ({cos(17*pi/10 r)},{sin(17*pi/10 r)});
      \coordinate (F) at ({cos(21*pi/10 r)},{sin(21*pi/10 r)});
      %線引き
      \draw[line width = 2pt] (F)--(B)--(C)--(D)--(E)--cycle;
      %nodeのスタイル指定
      \tikzset{block/.style={circle,white,fill=cyan,minimum size=.85cm,inner sep=0pt}};
      %node（テキストボックス的な）を指定座標に配置
      \node[inner sep=0pt] at (A) {\includegraphics[width=1cm]{pic/Logo.png}};
      \node[block] at (B) {\scriptsize MeCafe};
      \node[block] at (C) {\scriptsize メカ研};
      \node[block,fill=LightUIRed] at (D) {\scriptsize 電研};
      \node[block] at (E) {\scriptsize シス研};
      \node[block] at (F) {\scriptsize 情研};
    \end{tikzpicture}
  \end{textblock*}
\end{frame}

\begin{frame}{TechRingを創る仲間たち その４}{システム開発研究会}
  システム開発研究会にフォーカス\\
  活動の紹介
  \begin{textblock*}{0.4\linewidth}(280pt,180pt)
    \begin{tikzpicture}[scale=1]
      %座標設定
      \coordinate (A) at (0,-.1);
      \coordinate (B) at (0,1);
      \coordinate (C) at ({cos(9*pi/10 r)},{sin(9*pi/10 r)});
      \coordinate (D) at ({cos(13*pi/10 r)},{sin(13*pi/10 r)});
      \coordinate (E) at ({cos(17*pi/10 r)},{sin(17*pi/10 r)});
      \coordinate (F) at ({cos(21*pi/10 r)},{sin(21*pi/10 r)});
      %線引き
      \draw[line width = 2pt] (F)--(B)--(C)--(D)--(E)--cycle;
      %nodeのスタイル指定
      \tikzset{block/.style={circle,white,fill=cyan,minimum size=.85cm,inner sep=0pt}};
      %node（テキストボックス的な）を指定座標に配置
      \node[inner sep=0pt] at (A) {\includegraphics[width=1cm]{pic/Logo.png}};
      \node[block] at (B) {\scriptsize MeCafe};
      \node[block] at (C) {\scriptsize メカ研};
      \node[block] at (D) {\scriptsize 電研};
      \node[block,fill=LightUIRed] at (E) {\scriptsize シス研};
      \node[block] at (F) {\scriptsize 情研};
    \end{tikzpicture}
  \end{textblock*}
\end{frame}

\begin{frame}{TechRingを創る仲間たち その５}{情報処理研究会}
  情報処理研究会にフォーカス\\
  活動の紹介と出展の紹介を
  \begin{textblock*}{0.4\linewidth}(280pt,180pt)
    \begin{tikzpicture}[scale=1]
      %座標設定
      \coordinate (A) at (0,-.1);
      \coordinate (B) at (0,1);
      \coordinate (C) at ({cos(9*pi/10 r)},{sin(9*pi/10 r)});
      \coordinate (D) at ({cos(13*pi/10 r)},{sin(13*pi/10 r)});
      \coordinate (E) at ({cos(17*pi/10 r)},{sin(17*pi/10 r)});
      \coordinate (F) at ({cos(21*pi/10 r)},{sin(21*pi/10 r)});
      %線引き
      \draw[line width = 2pt] (F)--(B)--(C)--(D)--(E)--cycle;
      %nodeのスタイル指定
      \tikzset{block/.style={circle,white,fill=cyan,minimum size=.85cm,inner sep=0pt}};
      %node（テキストボックス的な）を指定座標に配置
      \node[inner sep=0pt] at (A) {\includegraphics[width=1cm]{pic/Logo.png}};
      \node[block] at (B) {\scriptsize MeCafe};
      \node[block] at (C) {\scriptsize メカ研};
      \node[block] at (D) {\scriptsize 電研};
      \node[block] at (E) {\scriptsize シス研};
      \node[block,fill=LightUIRed] at (F) {\scriptsize 情研};
    \end{tikzpicture}
  \end{textblock*}
\end{frame}

\begin{frame}{TechRingのこれから}{アッと驚く技術を、もっと広く世界へ}
  \begin{block}{大高専時代へ}
    \vspace{2mm}
    \begin{description}
      \setlength{\itemsep}{2mm}
      \item[創造の意欲] 学内イベントの開催により\\
      創作活動のモチベーションを上げ、活発化を目指す
      \item[幅広い視野] 学外への技術発信を行い\\
      アウトプットからフィードバックを得る
      \item[自立と友愛] チーム開発の経験から\\
      仲間とのコミュニケーションの重要さを学ぶ
    \end{description}
  \end{block}

  \vfill

  \begin{Large}
    モノを作るだけではなく \vfill \hspace{7em}\alert{コトを起こす}ものづくりを
  \end{Large}
\end{frame}

\begin{frame}[plain,noframenumbering]
  \begin{center}
    \begin{Large}
      ご清聴、ありがとうございました。
    \end{Large}
  \end{center}
\end{frame}

\end{document}